\subsection{Preliminaries}
\label{subsec:prelim}
% REVISED (2026-10-24): all definitions and theorems are now stated exactly as in the
% course slides (only the notation is adapted); explanations are outside the
% environments. The previous version is saved in /tmp/bk_1024 on the user's computer.

This subsection collects the notions on graphs and non-negative matrices used for both
models, following the course slides~\cite{baggio_graph,baggio_matrix,baggio_dt}.

\subsubsection{Graphs}
\label{subsubsec:prelim-graphs}

A \emph{weighted digraph} $G = (V, E, \{a_{ij}\}_{(i,j) \in E})$ has node set
$V = \{1, 2, \dots, n\}$, edge set $E \subseteq V \times V$ and a weight $a_{ij} > 0$ for
each edge $(i,j) \in E$; an edge $(i,i)$ is a \emph{self-loop}. Throughout the report
the nodes are agents, and an edge $(i,j)$ means that agent $i$ is influenced by agent
$j$. This is the convention of the course slides~\cite{baggio_dt}; the tutorial by
Proskurnikov and Tempo draws the edges in the opposite
direction~\cite{proskurnikov2017tutorial}.

\begin{definition}[Weighted adjacency matrix~\cite{baggio_graph}]
\label{def:adjacency}
Given a weighted digraph $G = (V, E,$ $\{a_{ij}\}_{(i,j) \in E})$ of order $n$ the
(weighted) adjacency matrix of $G$ is the $n \times n$ matrix $A$ in which:
\begin{itemize}
  \item[(i)] for each edge $(i,j) \in E$ the $(i,j)$ entry of $A$ is equal to the edge
             weight $a_{ij}$,
  \item[(ii)] all other entries of $A$ are equal to zero.
\end{itemize}
\end{definition}

\begin{definition}[Connectivity of digraphs~\cite{baggio_graph}]
A directed graph $G = (V, E)$ is
\begin{itemize}
  \item[(i)] \emph{strongly connected} if there is a directed walk from any node to any
             other node;
  \item[(ii)] \emph{weakly connected} if the undirected version of the graph is
              connected.
\end{itemize}
Moreover, $G$ has a
\begin{itemize}
  \item[(iii)] \emph{globally reachable node} if one of its nodes can be reached from
               any other node through a directed walk;
  \item[(iv)] \emph{directed spanning tree} if one of its nodes is the root of directed
              walks to every other node.
\end{itemize}
\end{definition}

With our convention, a directed walk from $i$ to $j$ means that agent $i$ is
influenced by agent $j$, either directly or through intermediaries. Hence $G$ is
strongly connected when every agent influences every other agent, at least
indirectly, and a globally reachable node is an agent that influences the whole
group.

\begin{definition}[Periodic digraph~\cite{baggio_graph}]
\label{def:periodic}
A strongly connected digraph $G = (V, E)$ is \emph{periodic} if there exists an
integer $k > 1$ (the period) that divides the length of every simple cycle of the
digraph.
\end{definition}

A strongly connected digraph that is not periodic is called \emph{aperiodic}.
Equivalently, a strongly connected digraph $G$ is periodic if and only if the greatest
common divisor of the lengths of its simple cycles is larger than
$1$~\cite{baggio_graph}. In particular,
a strongly connected digraph with a self-loop is aperiodic, since the self-loop is a
cycle of length $1$.

\begin{definition}[Strongly connected component~\cite{baggio_dt}]
Given a (weighted) digraph $G$. A subgraph $S$ is a \emph{strongly connected
component} of $G$ if $S$ is strongly connected and any other subgraph of $G$ strictly
containing $S$ is not strongly connected.
\end{definition}

\begin{definition}[Condensation digraph~\cite{baggio_dt}]
The \emph{condensation digraph} $C(G)$ of a (weighted) digraph $G$ is obtained by
collapsing each strongly connected component of $G$ into a single node. For any two
distinct strongly connected components $S_k$ and $S_h$ there is an edge $(S_k, S_h)$
in $C(G)$ iff there exists an edge $(i,j)$ in $G$ with $i \in S_k$ and $j \in S_h$.
\end{definition}

Self-loops are not allowed in $C(G)$~\cite{baggio_dt}. A node with no outgoing edges is
called a \emph{sink}~\cite{baggio_graph}.
A sink of $C(G)$ is a group of agents that are influenced only by members of the group
itself: no agent outside the group can change their opinions.

If $G$ contains a globally reachable node, then $C(G)$ has a unique
sink~\cite{baggio_dt}.

\subsubsection{Matrices}
\label{subsubsec:prelim-matrices}

\begin{definition}[Non-negative and stochastic matrices~\cite{baggio_matrix}]
\label{def:row-stochastic}
A matrix $A = [a_{ij}] \in \mathbb{R}^{n \times n}$ is
\begin{itemize}
  \item[(i)] \emph{non-negative} (\emph{positive}) if $a_{ij} \geq 0$ ($a_{ij} > 0$,
             resp.) for all $i, j \in \{1, \dots, n\}$;
  \item[(ii)] \emph{row-stochastic} if it is non-negative and
              $A \mathbf{1}_n = \mathbf{1}_n$;
  \item[(iii)] \emph{column-stochastic} if it is non-negative and
               $A^\top \mathbf{1}_n = \mathbf{1}_n$;
  \item[(iv)] \emph{doubly-stochastic} if it is row- and column-stochastic.
\end{itemize}
\end{definition}

\begin{theorem}[Spectral properties of row-stochastic matrices~\cite{baggio_matrix}]
\label{thm:stoch-spectrum}
For a row-stochastic matrix $A \in \mathbb{R}^{n \times n}$
\begin{itemize}
  \item[(i)] $1 \in \operatorname{spec}(A)$;
  \item[(ii)] $\operatorname{spec}(A)$ is a subset of the unit disk and $\rho(A) = 1$.
\end{itemize}
\end{theorem}

\begin{definition}[Semi-convergent and convergent matrices~\cite{baggio_matrix}]
\label{def:semiconv}
A matrix $A \in \mathbb{R}^{n \times n}$ is
\begin{itemize}
  \item[(i)] \emph{semi-convergent} if $\lim_{k \to +\infty} A^k$ exists (finite);
  \item[(ii)] \emph{convergent} (or \emph{Schur stable}) if it is semi-convergent and
              $\lim_{k \to +\infty} A^k = 0_{n \times n}$.
\end{itemize}
\end{definition}

\begin{theorem}[Convergence and spectral radius~\cite{baggio_matrix}]
\label{thm:spectral-radius}
Consider a matrix $A \in \mathbb{R}^{n \times n}$ and let
$\rho(A) := \max\{|\lambda| \mid \lambda \in \operatorname{spec}(A)\}$ (spectral
radius of $A$). Then:
\begin{itemize}
  \item[(i)] $A$ convergent $\iff$ $\rho(A) < 1$;
  \item[(ii)] $A$ semi-convergent and not convergent $\iff$ $1$ is a semisimple
              eigenvalue of $A$ and all the other eigenvalues have modulus strictly
              less than $1$.
\end{itemize}
\end{theorem}

\begin{definition}[Regularity~\cite{proskurnikov2017tutorial}]
\label{def:regular}
We call the matrix $A$ \emph{regular} if the limit
$A^\infty = \lim_{k \to \infty} A^k$ exists and \emph{fully regular} if, additionally,
the rows of $A^\infty$ are identical (that is, $A^\infty = \mathbf{1}_n w^\top$ for some
$w \in \mathbb{R}^n$).
\end{definition}

In the terminology of the course slides, a matrix is regular if and only if it is
semi-convergent (Definition~\ref{def:semiconv}).
% ADDED (2026-10-24): Definition 15 of the tutorial (notation W -> A, p_inf -> w),
% moved here from the FJ chapter because it is a general notion on matrices.

The properties of a non-negative matrix that determine its spectrum are characterized
by the graph notions above.

\begin{definition}[Irreducible and primitive matrices~\cite{baggio_matrix}]
A non-negative matrix $A \in \mathbb{R}^{n \times n}$, $n \geq 2$, is
\begin{itemize}
  \item[(i)] \emph{irreducible} if $\sum_{k=0}^{n-1} A^k$ is a positive matrix;
  \item[(ii)] \emph{primitive} if there exists $k \in \mathbb{N}$ such that $A^k$ is
              positive;
  \item[(iii)] \emph{reducible} if it is not irreducible.
\end{itemize}
\end{definition}

The next theorem uses two further notions~\cite{baggio_graph}: $\{I, J\}$ is a
\emph{partition} of $V = \{1, \dots, n\}$ if $I, J \neq \emptyset$,
$I \cap J = \emptyset$ and $I \cup J = V$; $P \in \{0,1\}^{n \times n}$ is a
\emph{permutation matrix} if it has exactly one entry equal to $1$ in every row and
every column.

\begin{theorem}[Irreducible matrices and digraphs~\cite{baggio_graph}]
\label{thm:irreducible}
Let $G = (V, E, \{a_{ij}\}_{(i,j) \in E})$ be a weighted digraph of order $n \geq 2$
with adjacency matrix $A$. The following statements are equivalent:
\begin{itemize}
  \item[(i)] $A$ is irreducible;
  \item[(ii)] there exists no permutation matrix $P$ such that $PAP^\top$ is block
              upper triangular;
  \item[(iii)] $G$ is strongly connected;
  \item[(iv)] for all partitions $\{I, J\}$ of $V$, there exists $i \in I$ and
              $j \in J$ such that $(i,j) \in E$.
\end{itemize}
\end{theorem}

\begin{theorem}[Primitive matrices and digraphs~\cite{baggio_graph}]
\label{thm:primitive}
Let $G = (V, E, \{a_{ij}\}_{(i,j) \in E})$ be a weighted digraph of order $n \geq 2$
with adjacency matrix $A$. The following statements are equivalent:
\begin{itemize}
  \item[(i)] $A$ is primitive,
  \item[(ii)] $G$ is strongly connected and aperiodic.
\end{itemize}
\end{theorem}

\begin{corollary}[Primitive matrices and digraphs with self-loops~\cite{baggio_graph}]
\label{cor:primitive-selfloop}
Let $G = (V, E,$ $\{a_{ij}\}_{(i,j) \in E})$ be a weighted digraph of order $n \geq 2$
with adjacency matrix $A$. If $G$ has at least one self-loop and is strongly
connected, then $A$ is primitive.
\end{corollary}

% REMOVED (to discuss): "If G is strongly connected, then G is aperiodic if and only
% if the period of each node is 1." -- the period of a node is not defined in this
% section, and the gcd characterization above is enough for what follows.

Given a weighted digraph $G$ with weighted adjacency matrix $A$, the \emph{weighted
out-degree} of node $i$ is the $i$-th row sum $\sum_{j=1}^{n} a_{ij}$ of
$A$~\cite{baggio_graph}. A row-stochastic matrix has all weighted out-degrees equal to
$1$; when some of them are smaller than $1$, the following notions apply.

\begin{definition}[Row-substochastic matrix~\cite{baggio_dt}]
\label{def:row-substochastic}
A non-negative matrix $A \in \mathbb{R}^{n \times n}$ is \emph{row-substochastic} if
its row-sums are at most $1$ and at least one row-sum is strictly less than $1$. That
is:
\[
  A \mathbf{1}_n \leq \mathbf{1}_n \quad \text{and} \quad
  \exists\, i \in \{1, \dots, n\} \text{ s.t. } e_i^\top A \mathbf{1}_n < 1.
\]
\end{definition}

\begin{lemma}[Convergent row-substochastic matrices~\cite{baggio_dt}]
\label{lem:substochastic}
Let $G$ be a weighted digraph with row-substochastic adjacency matrix $A$. Then, $A$ is
convergent if and only if $G$ contains directed walks from each node with (weighted)
out-degree $1$ to a node with (weighted) out-degree strictly less than $1$.
\end{lemma}
